\begin{table*}[ht]
\centering
\small
\begin{tabular}{lrrrrrrr}
\hline
& \multicolumn{4}{c}{\textbf{Python verifier (\%)}} & \multicolumn{3}{c}{\textbf{Lean verifier (\%)}} \\
\cmidrule(lr){2-5}\cmidrule(lr){6-8}
\textbf{Dataset} & True & False & Synth.\ fail & Exec.\ fail & True & Autoform.\ fail & Prover fail \\
\hline
GSM8K & 88.3 & 5.2 & 5.4 & 1.2 & 83.5 & 2.5 & 14.0 \\
MATH & 38.1 & 4.2 & 54.5 & 3.2 & 51.4 & 7.8 & 40.8 \\
OlympiadBench & 28.3 & 2.5 & 65.2 & 4.1 & 37.2 & 12.3 & 50.5 \\
OmniMATH & 24.7 & 3.3 & 69.8 & 2.1 & 29.7 & 14.5 & 55.9 \\
\hline
All & 35.5 & 3.4 & 58.1 & 3.0 & 43.2 & 10.8 & 46.0 \\
\hline
\end{tabular}
\caption{Distribution of verifier outcomes on the annotated ProcessBench steps. For Python, \emph{Synth.\ fail} means no executable checker was produced (including programs rejected by the sandbox) and \emph{Exec.\ fail} means the checker did not return a Boolean (runtime error, time-out). For Lean, \emph{Autoform.\ fail} means no well-formed theorem was produced and \emph{Prover fail} that no candidate proof compiled; the Lean pipeline never returns an explicit False, so every non-True outcome counts as an incorrect-step verdict.}
\label{tab:verifier_outcomes}
\end{table*}
